\section{Contesto e obiettivi}

\subsection{Descrizione del problema}
Il Laboratorio in Campo coinvolge gli studenti del secondo anno di Scienze Ambientali in escursioni
guidate dai docenti. Gli studenti, divisi in gruppi, devono raccogliere osservazioni sul territorio:
posizione, descrizione, misure, fotografie e schede compilate.
\subsection{Stakeholder}
\begin{center}
\begin{tblr}{width=\linewidth, colspec={X[l]X[l]X[l]}, cells={font=\small},
             row{1}={bg=arancio}, hline{2-Z}={wd=0.5pt, fg=arancio}}
Stakeholder & Ruolo & Interesse principale \\
Studenti (secondo anno) & Utenti finali & Raccogliere e consegnare i dati dell'uscita in modo semplice \\
Docenti del laboratorio (12) & Committenti e utenti & Definire le attività e valutare i dati raccolti \\
Ateneo (Ca' Foscari) & Ente promotore & Qualità didattica e rispetto delle norme sui dati \\
\end{tblr}
\end{center}

\subsection{Ambito e vincoli}
L'app è per Android. Le uscite avvengono in ambiente montano, dove la copertura di rete può mancare.
